\documentclass[11pt,a4paper]{article}
\usepackage[margin=2.4cm]{geometry}
\usepackage{amsmath,amssymb,amsthm,mathtools}
\usepackage{booktabs}
\usepackage{enumitem}
\usepackage{xcolor}
\usepackage{hyperref}

\hypersetup{
  colorlinks=true,
  linkcolor=blue!70!black,
  citecolor=blue!70!black,
  urlcolor=blue!70!black,
  pdftitle={Cuboctahedral Rotator Calculus, Part II: Post Functional Completeness, 3-SAT Phase-Polynomial Reduction, Constructive de Rham Cohomology, and Finite-State Adversary Barriers on Z_512 in Lean 4},
  pdfauthor={Dmitrii A. Karataev}
}

\theoremstyle{plain}
\newtheorem{theorem}{Theorem}[section]
\newtheorem{lemma}[theorem]{Lemma}
\newtheorem{proposition}[theorem]{Proposition}
\newtheorem{corollary}[theorem]{Corollary}
\theoremstyle{definition}
\newtheorem{definition}[theorem]{Definition}
\newtheorem{remark}[theorem]{Remark}

\title{\textbf{Cuboctahedral Rotator Calculus, Part~II:\\ Post Functional Completeness, 3-SAT Phase-Polynomial Reduction, Constructive de~Rham Cohomology, and Finite-State Adversary Barriers on $\mathbb{Z}_{512}$ in Lean~4}}
\author{Dmitrii A.~Karataev}
\date{}

\begin{document}
\maketitle

\begin{abstract}
In Part~II of this series, we establish the Boolean logic, algebraic complexity, cochain cohomology, and measure-theoretic foundations of the cuboctahedral rotator calculus on $\mathbb{Z}_{N_{\mathrm{ring}}(d)}$ ($N_{\mathrm{ring}}(3)=512$, $Q_\square(3)=128$, $\Phi_{\mathrm{half}}(3)=256$) and the $12$-vertex cuboctahedral complex $O_2(3)$. We prove that on the antipodal pole subgroup $\{0,256\}\subset\mathbb{Z}_{512}$, pre-rotator phase addition (Channel~1, Post's linear class~$L$: $\{\mathrm{NOT},\mathrm{XOR}_3\}$) and post-rotator projection summation with integer division by $512$ (Channel~2, non-linear threshold class: $\{\mathrm{OR}_3,\mathrm{AND}_3\}$) form a functionally complete Boolean system by Post's theorem, reducing $M$-clause 3-SAT satisfiability to $\sum_{m=1}^M q_{\mathrm{OR}}(\psi_{m,1},\psi_{m,2},\psi_{m,3}) = 0$. We prove the $7\to 1$ remainder-gluing law of forward verification, identical annihilation of the Channel~2 linear gradient on balanced $4$-clause Channel~1 parity loops ($\mathrm{xor4LinearGradientSum}\equiv 0$), and $7^T$ ray-preimage branching across the depth-$T$ cuboctahedral coordination ball $\mathcal{V}_{O_2}(3,T) = (10T^3+15T^2+11T+3)/3$.

Crucially, we analyze why cryptographic security against any \emph{realizable} adversary depends not on an abstract external Turing machine with an infinite tape outside the system, but on finite-state observers embedded \emph{inside} the state-transition graph: (1)~any embedded sub-automaton obeys four formal information-capacity and causal-horizon barriers (Dirichlet self-inclusion capacity, $72^2 = 5\,184$ projection blindness from the internal $M$-carrier to the external $T$-observable, cone-crossing latency, and coboundary back-action); and (2)~because external $T$-layer observation extracts only a $1$-bit sign from the $72$-bit four-channel $\mathbb{Z}_{512}$ Verlet state cell while leaving $71$ internal phase bits algebraically unobservable without breaking the $512\mathbb{Z}$-quantized plaquette holonomy, the $M\to T$ channel embeds a native Shannon one-time pad (Vernam cipher, $H(M\mid T)=H(M)$) that is information-theoretically immune to complexity assumptions ($P$ vs.\ $NP$). Finally, we prove the constructive de~Rham--Stokes--Gauss--Chern cochain identities ($d_1\circ d_0\equiv 0$, $d_2\circ d_1\equiv 0$), division-free Neumann nilpotent inversion ($N_{\mathrm{nil}}^3=0$), and Banach--Tarski-free finite measure additivity in pure \textbf{Lean~4} (\texttt{thm4}, \texttt{thm5}, \texttt{0 sorryAx}, \texttt{0 Classical.choice}).

\medskip\noindent
\textbf{Keywords:} Post functional completeness, 3-SAT reduction, internal observer barriers, Shannon Vernam one-time pad, discrete de~Rham cohomology, nilpotent jets, constructive measure, Lean~4.

\smallskip\noindent
\textbf{MSC 2020:} 03B05, 68Q17, 94A60, 55N35, 28A05, 68V20.
\end{abstract}

\section{Post Functional Completeness, Branchless 3-SAT Representation on \texorpdfstring{$\mathbb{Z}_{512}$}{Z512}, and Algebraic Asymmetry of Verification vs.\ Inversion}
\label{sec:boolean-3sat}

On the antipodal poles $\{0,\Phi_{\mathrm{half}}(d)\}=\{0,256\}\subset\mathbb{Z}_{N_{\mathrm{ring}}(d)}$ ($\mathcal{C}_3(0)=+128$ encoding truth~$1$, $\mathcal{C}_3(256)=-128$ encoding falsity~$0$), the piecewise-linear rotator induces two canonical algebraic channels on triangular faces $(\psi_1,\psi_2,\psi_3)\in\{0,256\}^3$:
\begin{enumerate}[leftmargin=*,noitemsep]
\item \textbf{Channel~1 (Post's affine linear class~$L$ --- pre-rotator phase addition):}
  the half-turn shift $\mathrm{NOT}(\psi)=(\psi+256)\bmod 512$ negates the projection ($\mathcal{C}_3(\psi+256)=-\mathcal{C}_3(\psi)$), while phase addition yields three-input parity $\mathrm{XOR}_3$:
  \begin{equation}
    \label{eq:pre-rotator-xor}
    Q_\square(d)^2\,\mathcal{C}_d(\psi_1+\psi_2+\psi_3) \;\equiv\; \mathcal{C}_d(\psi_1)\,\mathcal{C}_d(\psi_2)\,\mathcal{C}_d(\psi_3),
  \end{equation}
  where the pole subgroup $(\{0,256\},+\bmod 512)\cong(\mathbb{F}_2,\oplus)$ ($256+256\equiv 0\pmod{512}$).
\item \textbf{Channel~2 (Non-linear threshold class --- post-rotator projection sum with integer division by $N_{\mathrm{ring}}(d)=512$):}
  set $\Delta_{\mathrm{3SAT}}(\psi_1,\psi_2,\psi_3) := Q_\square(d)-\bigl(\mathcal{C}_d(\psi_1)+\mathcal{C}_d(\psi_2)+\mathcal{C}_d(\psi_3)\bigr)$. Since $(1+d)Q_\square(d)=4\times 128=512=N_{\mathrm{ring}}(d)$ at $d=3$, truncated integer division by $N_{\mathrm{ring}}(d)$ defines branchless disjunction and conjunction:
  \begin{equation}
    \label{eq:post-rotator-or-and}
    q_{\mathrm{OR}}(\psi_1,\psi_2,\psi_3) := \mathrm{tdiv}\bigl(\Delta_{\mathrm{3SAT}}(\psi_1,\psi_2,\psi_3),\,N_{\mathrm{ring}}(d)\bigr)\in\{0,1\},\qquad
    q_{\mathrm{AND}}(\psi_1,\psi_2,\psi_3) := \mathrm{tdiv}\Bigl(Q_\square(d)+\sum_{k=1}^3\mathcal{C}_d(\psi_k),\,N_{\mathrm{ring}}(d)\Bigr)\in\{0,1\}.
  \end{equation}
\end{enumerate}

\begin{theorem}[Post Functional Completeness \{\texttt{NOT}, \texttt{XOR}, \texttt{OR}, \texttt{AND}\}, 3-SAT Representation, $7\to 1$ Remainder Gluing, and Two-Channel Non-Commutativity on $\mathbb{Z}_{512}$]
\label{thm:main-4}
For $d=3$, the following exact relations hold across the pole octants $\{0,256\}^3\subset\mathbb{Z}_{512}^3$ (\texttt{thm4\_boolean\_post\_completeness\_3sat\_reduction\_and\_channel\_orthogonality}):
\begin{enumerate}[label=\textbf{(\roman*)},leftmargin=*]
\item \textbf{Post functional completeness and equivalence of 3-SAT satisfiability to zero Channel~2 quotients:}
  across all $2^3=8$ states in $\{0,256\}^3$, the multiplicative parity identity~\eqref{eq:pre-rotator-xor} holds, $q_{\mathrm{AND}}$ equals~$1$ uniquely at $(0,0,0)$ and~$0$ on the other $7$~octants, while the defect quotient $q_{\mathrm{OR}}$ equals~$0$ on all $7$~satisfying octants of a 3-SAT clause $(l_1\lor l_2\lor l_3)$ and equals~$1$ ($\Delta_{\mathrm{3SAT}}=512$) uniquely on the falsifying octant $(256,256,256)$. Consequently, by the polynomial equivalence between Turing machines and cellular automata ($\mathbf{P}_{\mathrm{TM}}\equiv\mathbf{P}_{\mathrm{CA}}$, Cook--Levin), every Boolean circuit decomposes into Channel~1 and Channel~2, and a 3-SAT formula of $M$ clauses is satisfiable if and only if $\sum_{m=1}^M q_{\mathrm{OR}}(\psi_{m,1},\psi_{m,2},\psi_{m,3}) = 0$.
\item \textbf{Non-injective $7\to 1$ remainder gluing in forward verification ($\mathbf{NP}$):}
  the forward quotient operator $q_{\mathrm{OR}}=\mathrm{tdiv}(\Delta_{\mathrm{3SAT}},512)$ discards the three-valued remainder $r\in\{-256,0,+256\}$ and maps $\mathrm{count3SatSatisfyingOctants}(3)=2^3-1=7$ distinct input octants onto the single output $q_{\mathrm{OR}}=0$ ($1\to 7$ preimage branching per clause under inversion).
\item \textbf{Identical annihilation of the Channel~2 linear gradient on Channel~1 parity constraints and channel non-commutativity:}
  across the four 3-SAT clauses encoding a single Channel~1 parity equation ($x_1\oplus x_2\oplus x_3=0$), the sum of linear projections $\mathcal{C}_3$ vanishes identically on all $8$~octants ($\mathrm{xor4LinearGradientSum}(3,\psi_1,\psi_2,\psi_3)\equiv 0$), while the rotator $\mathcal{C}_3$ and the Channel~2 threshold quotient $q_{\mathrm{AND}}$ are non-linear and non-distributive with respect to Channel~1 phase addition.
\item \textbf{Single-point depth-$T$ inverse cone underdetermination:}
  when reconstructing a state at depth $T$ steps from a single observation node $\mathbf{x}_0$, the $T$-th coordination shell of the cuboctahedron contains $\mathrm{cuboctShellNodes}(3,T)=10T^2+2$ equidistant nodes, the cumulative cone volume $\mathcal{V}_{O_2}(3,T)=\mathrm{cuboctBallNodes}(3,T)$ obeys the exact cubic law $3\,\mathcal{V}_{O_2}(3,T)=10T^3+15T^2+11T+3$ ($\mathcal{V}_{O_2}(3,3)=3\cdot 7^2=147$), the Channel~2 gluing $7\to 1$ induces $\mathrm{preimageRayBranching}(3,T)=7^T$ indistinguishable preimages along each ray ($7^3=343$ at $T=3$), and balanced 4-clause parity loops ($\mathrm{xor4LinearGradientSum}\equiv 0$) form zero-gradient cloaked kernels invisible to linear probing at $\mathbf{x}_0$.
\end{enumerate}
\end{theorem}

\section{Embedded Finite-State Adversaries vs.\ External Turing Machines and the Built-In Shannon Vernam Cipher on \texorpdfstring{$\mathbb{Z}_{512}$}{Z512}}
\label{sec:embedded-adversary-vernam}

In the model-theoretic framework of A.\,I.~Mal'cev (\emph{Algebraic Systems}~\cite{Malcev1973}), the calculus is a two-sorted algebraic structure $\mathfrak{S} = \langle \mathcal{M},\,\mathcal{T};\,\pi_{M\to T},\,\iota_{T\to M}\rangle$, where the \emph{carrier} (\emph{носитель}) of a sort is its underlying domain set: the \textbf{discrete $M$-carrier} $\mathcal{M} = \mathbb{Z}_{512}$ (internal modular phase states on the cuboctahedral lattice) and the \textbf{continuous $T$-carrier} $\mathcal{T}$ (externally observable Euclidean/threshold invariants).

Classical complexity-based cryptography models an adversary as an abstract Turing machine placed \emph{outside} the state space with an unbounded external tape and zero signal-propagation cost, leading to the fear that either a polynomial-time 3-SAT algorithm ($P=NP$) or quantum speedup could collapse cryptographic security. Theorem~\ref{thm:main-4} and Theorem~\ref{thm:main-5} expose two structural facts that decouple realizable channel security from abstract external-tape assumptions:
\begin{enumerate}[label=\textbf{(\arabic*)},leftmargin=*]
\item \textbf{Any realizable adversary is an embedded finite sub-automaton subject to four formal inversion barriers:}
  an actual cryptanalytic solver is not an oracle outside the universe, but a finite sub-complex occupying $V_{\mathrm{adv}} < \mathcal{V}_{O_2}(3,T)$ nodes within the same causal cone. Consequently, any embedded solver is bounded by:
  \begin{itemize}[noitemsep]
  \item \emph{Barrier~I (Dirichlet self-inclusion register bound):} a sub-complex of $V_{\mathrm{adv}}$ cells cannot store the $B_2\cdot\mathcal{V}_{O_2}(3,T)$ state bits of a strictly larger causal cone ($B_2 V_{\mathrm{adv}} < B_2\mathcal{V}_{O_2}(3,T)$, where $B_2 = 8\times 9 = 72$ bits per four-channel symplectic cell);
  \item \emph{Barrier~II (Projection blindness square law $K_P(3)\cdot 81 = 72^2 = 5\,184$):} a single coarse $T$-layer threshold observation extracts only $1$~bit ($q_{\mathrm{OR}}\in\{0,1\}$) out of the $72$~internal phase and momentum bits of the $M$-carrier cell, leaving $7\to 1$ octant gluing per clause and $7^T$ ray-preimage branching at depth~$T$;
  \item \emph{Barrier~III (Cone-horizon propagation latency):} gathering boundary data across radius $T$ requires at least $T$ steps, by which time the causal horizon has expanded to $2T > T$;
  \item \emph{Barrier~IV (Coboundary back-action):} writing an inferred state into the solver's own registers alters the local $1$-cochain connection $r = \mathrm{edgeConn}(3,\Delta\varphi)$ and shifts the $512\,n_f$ plaquette coboundary.
  \end{itemize}
\item \textbf{Built-in information-theoretic Vernam one-time pad on the discrete $M$-carrier $\mathbb{Z}_{512}$:}
  because external $T$-layer observation accesses only the coarse $1$-bit quotient while the remaining $71$~internal phase bits per cell (and the $r\in\{-256,0,+256\}$ remainder across zero-gradient parity kernels $\mathrm{xor4LinearGradientSum}\equiv 0$) remain algebraically orthogonal to the $T$-layer projection and topologically protected by the integer Chern coboundary $(d_1 r)_f = 512\,n_f$, phase-locked transmission on the $M$-carrier acts as a native Shannon Vernam one-time pad ($H(M\mid T) = H(M)$). By Shannon's perfect-secrecy theorem, an information-theoretic Vernam pad cannot be broken by any computational power whatsoever---even if $P=NP$ held for abstract external Turing machines.
\end{enumerate}

\section{Constructive de~Rham--Stokes--Gauss--Chern Cohomology, Division-Free Linear Algebra, and Choice-Free Measure Additivity}
\label{sec:constructive-foundations}

\begin{theorem}[Constructive de~Rham--Stokes--Gauss--Chern Cohomology, Division-Free Nilpotent and Symplectic Inversion, and Banach--Tarski-Free Measure Additivity]
\label{thm:main-5}
For $d=3$ on $\mathbb{Z}_{512}$ and the cuboctahedral complex~$O_2(3)$, the following exact constructive identities hold (\texttt{thm5\_constructive\_de\_rham\_cohomology\_nilpotent\_algebra\_and\_banach\_tarski\_free\_measure}):
\begin{enumerate}[label=\textbf{(\roman*)},leftmargin=*]
\item \textbf{Euclidean connection splitting $\Delta\varphi = 512\,q + r$ and variational maximum on the $L^1$ diamond face:}
  for all $\Delta\varphi\in[-512,512]$, Euclidean floor division $q = \mathrm{branchQuot}(3,\Delta\varphi) = \lfloor(\Delta\varphi+255)/512\rfloor$, $r = \mathrm{edgeConn}(3,\Delta\varphi)\in(-256,256]$ satisfies $512\,q + r \equiv \Delta\varphi$; moreover, along $\varphi = 256 - k$ ($k\in[0,128]$), the conjugate projections $V(k) = |\mathcal{C}_3(256-k)| = 128-k$ and $I(k) = \mathcal{S}_3(256-k) = k$ obey $V(k)+I(k)\equiv 128 = Q_\square(3)$, while their quadratic product $P(k) = V(k)I(k) = (128-k)k$ attains its unique maximum $P_{\max} = K_P(3)^2 = 64^2 = 4\,096$ at $k^* = K_P(3) = 64$.
\item \textbf{Constructive Stokes theorem ($d_1\circ d_0\equiv 0$) and integer Chern class quantization:}
  for any $0$-cochain $\varphi = (\varphi_0,\varphi_1,\varphi_2,\varphi_3)\in\mathbb{Z}^4$ around a closed $4$-cycle face~$f$, the raw coboundary vanishes identically ($(d_1(d_0\varphi))_f \equiv 0$), so that the connection circulation $(d_1 r)_f = \sum_{e\in\partial f}\mathrm{edgeConn}(3,(d_0\varphi)_e)$ is strictly an integer multiple of the ring period:
  \[
    (d_1 r)_f \;=\; 512\,n_f \;\in\; 512\,\mathbb{Z},\qquad
    n_f := -\sum_{e\in\partial f} q_e \in \mathbb{Z}.
  \]
\item \textbf{Constructive Gauss--Ostrogradsky divergence theorem ($d_2\circ d_1\equiv 0$ on the closed cube surface $\partial\square^3$):}
  for any integer $1$-cochain $(e_1,\dots,e_{12})\in\mathbb{Z}^{12}$ on the $12$~edges of a cube, the sum of oriented face fluxes $(d_1 e)_f$ over all $|\square(3)|=6$ boundary faces vanishes identically ($\mathrm{cubeClosedSurfaceFluxSum}\equiv 0$).
\item \textbf{Division-free linear algebra ($\mathrm{SL}_2(\mathbb{Z})$ and nilpotent Neumann inversion):}
  the symplectic Verlet map $(u,v)\mapsto(2u-v+\Phi,u)$ ($\det = +1\in\mathrm{SL}_2(\mathbb{Z})$) inverts in place via $(u',v')\mapsto(v',2v'-u'+\Phi)$, while the 4D nilpotent operator $\mathcal{M} = I_4 + N_{\mathrm{nil}}$ ($N_{\mathrm{nil}}^3 = 0$, $\det(2\mathcal{M}) = 24$) inverts exactly via the finite 3-term Neumann polynomial $\mathcal{M}^{-1} = I_4 - N_{\mathrm{nil}} + N_{\mathrm{nil}}^2$.
\item \textbf{Choice-free, Banach--Tarski-free finite additivity of cuboctahedral measure:}
  partitioning the cuboctahedral ball $\mathcal{B}_{O_2}(3,T)$ into the central node and disjoint coordination shells $\mathrm{cuboctShellNodes}(3,k) = 10k^2+2$ strictly preserves integer volume $1 + \sum_{k=1}^T(10k^2+2) = \mathcal{V}_{O_2}(3,T)$, with $12\times(4^2+5^2) = 492 > 256 = \Phi_{\mathrm{half}}(3)$.
\end{enumerate}
\end{theorem}

\section{Formal Verification in Lean~4}
\label{sec:lean4}

The formalization of Theorems~\ref{thm:main-4} and~\ref{thm:main-5} (\texttt{thm4\_boolean\_post\_completeness\_3sat\_reduction\_and\_channel\_orthogonality} and \texttt{thm5\_constructive\_de\_rham\_cohomology\_nilpotent\_algebra\_and\_banach\_tarski\_free\_measure}) resides in \texttt{RotatorCalculusPaper.lean} in \url{https://github.com/KarataevDmitry/cuboctahedral-rotator-calculus} (\texttt{0 sorryAx}, \texttt{0 Classical.choice}).

\begin{thebibliography}{9}
\bibitem{Malcev1973}
A.\,I.~Mal'cev, \emph{Algebraic Systems}, Die Grundlehren der mathematischen Wissenschaften, vol.~192, Springer-Verlag, Berlin--Heidelberg--New York, 1973 (Russian original: \emph{Алгебраические системы}, Nauka, Moscow, 1970).
\bibitem{Post1941}
E.\,L.~Post, \emph{The Two-Valued Iterative Systems of Mathematical Logic}, Annals of Mathematics Studies, no.~5, Princeton University Press, 1941.
\bibitem{Shannon1949}
C.\,E.~Shannon, ``Communication theory of secrecy systems,'' \emph{Bell System Technical Journal} \textbf{28}(4) (1949), 656--715.
\bibitem{Moura2021}
L.~de~Moura and S.~Ullrich, ``The Lean~4 theorem prover and programming language,'' in \emph{CADE~28}, LNCS vol.~12699, Springer, 2021, pp.~625--635.
\end{thebibliography}

\end{document}
